% source: Proposition A.3, preparation cost laws; the nonzero-price posterior argument in O2_proof.tex.
% Manuscript location: appendix/A6_robustness.tex:103-113, Steps 1 and 2 of the proof of prop:atomless.
% The low component has support [0, bar c_0], and the high component has support [c_-,c_+].

\emph{Step 1: the floor at the nonzero orders.} Fix a candidate in either economy. A price reached by a
nonzero order has its posterior in \([m,M]\), so
\[
B_{F_i}(\bar\mu)\ge B_{F_i}(m)\ge B_1(m)>\bar c_0.
\]
Every cost in the support of \(\Phi^0\) therefore prepares strictly at both nonzero-order prices, giving entry
at least \(\rho\) after a purchase and after a sale.

\emph{Step 2: orders and preparation by the low component.} At \(F_0\), entry in \([0,1]\) gives correct gross
returns at most \(\Delta_0<k\), and wrong-signed orders earn at most \(-k\). Both classes therefore abstain.
Every order and price has posterior \(\pi\), and
\(\bar c_0<B_1(m)\le B_0(\pi)<c_-\) gives entry \(\rho\).
At \(F_1\), the nonzero-order floor and the trading bound give correct gross returns at least
\(\rho\kappa\Delta_1>k\), so both classes place full orders. Their posteriors are \((m,\pi,M)\), and each price
belief is an average of these. The low component therefore prepares strictly at every price. These order
comparisons hold under both investor tie rules. The price formula and the strict support windows give the
equilibrium tuples and outcomes of Proposition~\ref{prop:atomless}.
